% Generated draft: requires mathematical and source review.
\documentclass[11pt]{article}
\usepackage[margin=25mm]{geometry}
\usepackage{fontspec}
\setmainfont{Latin Modern Roman}
\usepackage{amsmath,amssymb,mathtools,unicode-math}
\setmathfont{Latin Modern Math}
\usepackage{enumitem,needspace,xurl,hyperref}
\setlength{\parindent}{0pt}
\setlength{\parskip}{5pt}
\setlength{\emergencystretch}{3em}
\newcommand{\N}{\mathbb{N}}
\newcommand{\Z}{\mathbb{Z}}
\newcommand{\Q}{\mathbb{Q}}
\newcommand{\R}{\mathbb{R}}
\newcommand{\C}{\mathbb{C}}
\newcommand{\F}{\mathbb{F}}
\newcommand{\A}{\mathbb{A}}
\newcommand{\PP}{\mathbb{P}}
\newcommand{\E}{\mathbb{E}}
\newcommand{\eps}{\varepsilon}
\newcommand{\dd}{\,\mathrm{d}}
\begin{document}
\section*{UnsolvedMath ID: 30002762}
\textbf{Source problem code:} catalogue entry
\subsection{Definitions and mathematical statement}
Let \(G\) be a finitely presented group, and let \(S\subseteq G\) be a finite set whose conjugates generate \(G\). Put
\[
T=\{h s^{\varepsilon}h^{-1}:h\in G,\ s\in S,\ \varepsilon\in\{-1,1\}\}.
\]
The associated conjugation-invariant word norm is
\[
\|x\|_S=\min\{k\in\Z_{\geq0}:x=t_1\cdots t_k,\ t_i\in T\},
\]
where the empty product is the identity. It defines the bi-invariant metric \(d_S(x,y)=\|x^{-1}y\|_S\).

For \(g\in G\), the cyclic subgroup \(\langle g\rangle\) is bounded if
\[
\sup_{n\in\Z}\|g^n\|_S<\infty.
\]
Define its stable length by
\[
\tau_S(g)=\lim_{n\to\infty}\frac{\|g^n\|_S}{n}
=\inf_{n\geq1}\frac{\|g^n\|_S}{n}.
\]
For an infinite cyclic subgroup, undistortion means that the orbit map \(n\mapsto g^n\), from \(\Z\) with its usual word metric to \((G,d_S)\), is a quasi-isometric embedding; equivalently, \(\tau_S(g)>0\). Every unbounded cyclic subgroup is infinite, so this convention suffices for the question.

A homogeneous quasimorphism is a map \(\phi:G\to\R\) such that
\[
\sup_{x,y\in G}|\phi(xy)-\phi(x)-\phi(y)|<\infty,
\qquad
\phi(x^n)=n\phi(x)\quad(x\in G,\ n\in\Z).
\]
Here, detection of \(\langle g\rangle\) means that some homogeneous quasimorphism \(\phi\) satisfies \(\phi(g)\ne0\).

\textbf{Question.} Do there exist a finitely presented group \(G\), a finite set \(S\) as above, and an element \(g\in G\) such that at least one of the following holds?
\begin{enumerate}
\item \(\displaystyle\sup_{n\in\Z}\|g^n\|_S=\infty\) and \(\tau_S(g)=0\).
\item \(\displaystyle\sup_{n\in\Z}\|g^n\|_S=\infty\) and \(\phi(g)=0\) for every homogeneous quasimorphism \(\phi:G\to\R\).
\end{enumerate}
These are respectively counterexamples to the bounded-or-undistorted and bounded-or-detected dichotomies. The supplied record does not explicitly define undistortion or detection; the standard interpretations used above require confirmation against the source. Symmetrization of the generating conjugacy classes is made explicit in the definition of \(T\).
\subsection{Short English statement}
Find a finitely presented group with a word norm generated by finitely many conjugacy classes that has an unbounded cyclic subgroup whose powers either grow sublinearly in that norm or are annihilated by every real-valued homogeneous quasimorphism. The supplied record labels this problem open; this draft does not independently verify that status.
\subsection{Sources}
\begin{itemize}
\item Local UnsolvedMath registry: supplied complete record for permanent ID \(30002762\), title ``Finitely Presented Counterexamples for Conjugation-Invariant Norms,'' source problem code \texttt{catalogue entry}, Oberwolfach Reports Open Problems collection, proposed year 2015. No individual proposer is supplied.
\item \emph{Geometric Topology} (2015), Oberwolfach Reports, pp. 187--233, DOI \texttt{10.4171/owr/2015/3}. Supplied source report and literature-source link, labeled ``original OWR report'': \url{https://doi.org/10.4171/owr/2015/3}.
\item Brandenbursky, Gal, Kedra, and Marcinkowski, supplied literature-source label ``bi-invariant word metrics,'' DOI \texttt{10.1017/S0017089515000129}: \url{https://doi.org/10.1017/S0017089515000129}. The supplied record does not give fuller bibliographic details.
\item Supplied dated literature curation, checked \texttt{2026-08-21}: status ``open.'' Its assessment states that the cited work establishes the two dichotomies for many geometric classes and gives counterexamples outside the finitely presented setting, but that no finitely presented group violating either dichotomy was identified. This is supplied dated curation derived from the OWR record, not an independently verified literature assessment; its time-sensitive status and cited claims require independent verification.
\item The registry labels the statement \texttt{corrected\_verified} and reports that it was checked against the official Oberwolfach report, preserving its hypotheses, notation, and open target. This is a supplied verification claim, not verification performed for this draft.
\end{itemize}
\end{document}
